% GENERATED FILE. Edit canonical lessons, then run npm run generate:books.
% canonical-source-hash: fnv1a64:9ddd8991cdca6bbf
% canonical-lessons: LA-C484-fagus, LA-C484-abies, LA-C484-ulmus, LA-C484-fraxinus, LA-C484-opacus

\chapter{Beech, Fir, Elm, Ash, Shady}
\label{ch:la-beech-fir-elm-ash-shady}
\hlchaptermodality{484}

\begin{chapteropening}
\textbf{By the end of this chapter:} \emph{I can name the trees of a forest and describe it: beech, fir, elm and ash, and shady woods.}

Complete the last lesson of chapter 484: I can name the trees of a forest and describe it: beech, fir, elm and ash, and shady woods.
\end{chapteropening}

\section[fāgus, fāgī]{fāgus, fāgī — a beech tree}

\label{lesson:LA-C484-fagus}



\begin{quote}
Before the new one: say the Latin for to trickle, then the Latin for to flood.
\end{quote}

\subsection*{You'll want to know: fāgus, fāgī}
\textbf{fāgus, fāgī} — \textquotedblleft{}a beech tree\textquotedblright{}.

\begin{quote}
Pastor sub fāgō sedet et cantat.
\end{quote}

\begin{quote}
\emph{The shepherd sits under a beech tree and sings.}
\end{quote}

Second declension but feminine, like most tree names: \emph{fāgus, fāgī f.} English \emph{beech} is a Germanic cousin, not a borrowing.

\begin{cousinweb}
\emph{fāgus} is inherited from Proto-Italic \emph{fāgos}, from PIE \emph{b\textsuperscript{h}eh\textsubscript{2}ǵos} 'beech tree' (EtymDB). English \emph{beech} and German \emph{Buche} are Germanic cousins (Lewis \& Short).
\end{cousinweb}

\begin{grammarlens}[title={What you've built}]
One more word for telling, explaining and describing at B1.
\end{grammarlens}

\subsection*{Your turn}
\begin{itemize}
\raggedright
  \item \emph{Say it:} \emph{fāgus, fāgī}
  \item \emph{Say it:} \emph{fāgus, fāgī}, once more
  \item \emph{Say it:} a sentence of your own with \emph{fāgus, fāgī}
  \item \emph{Recall:} say the Latin for to trickle, then the Latin for to flood, then say \emph{fāgus, fāgī} again
\end{itemize}

\begin{tcolorbox}[breakable,colback=teal!4,colframe=teal!35!black,title={Before you move on}]
What does fāgus, fāgī mean? (\textquotedblleft{}A beech tree\textquotedblright{}.) Say it once more.
\end{tcolorbox}
\section[abiēs, abietis]{abiēs, abietis — a fir tree; fir wood}

\label{lesson:LA-C484-abies}



\begin{quote}
Before the new one: say the Latin for to flood, then the Latin for a beech tree.
\end{quote}

\subsection*{You'll want to know: abiēs, abietis}
\textbf{abiēs, abietis} — \textquotedblleft{}a fir tree; fir wood\textquotedblright{}.

\begin{quote}
In monte altō abietēs crēscunt.
\end{quote}

\begin{quote}
\emph{Fir trees grow on the high mountain.}
\end{quote}

Third declension, feminine: \emph{abiēs, abietis f.} Its wood was light and good for ships and oars.

\begin{cousinweb}
\emph{abiēs} is inherited from Proto-Italic \emph{abiets} 'fir' (EtymDB); Lewis \& Short call its further origin uncertain.
\end{cousinweb}

\begin{grammarlens}[title={What you've built}]
One more word for telling, explaining and describing at B1.
\end{grammarlens}

\subsection*{Your turn}
\begin{itemize}
\raggedright
  \item \emph{Say it:} \emph{abiēs, abietis}
  \item \emph{Say it:} \emph{abiēs, abietis}, once more
  \item \emph{Say it:} a sentence of your own with \emph{abiēs, abietis}
  \item \emph{Recall:} say the Latin for to flood, then the Latin for a beech tree, then say \emph{abiēs, abietis} again
\end{itemize}

\begin{tcolorbox}[breakable,colback=teal!4,colframe=teal!35!black,title={Before you move on}]
What does abiēs, abietis mean? (\textquotedblleft{}A fir tree; fir wood\textquotedblright{}.) Say it once more.
\end{tcolorbox}
\section[ulmus, ulmī]{ulmus, ulmī — an elm tree}

\label{lesson:LA-C484-ulmus}



\begin{quote}
Before the new one: say the Latin for a beech tree, then the Latin for a fir tree.
\end{quote}

\subsection*{You'll want to know: ulmus, ulmī}
\textbf{ulmus, ulmī} — \textquotedblleft{}an elm tree\textquotedblright{}.

\begin{quote}
Agricolae vītēs ad ulmōs ligant.
\end{quote}

\begin{quote}
\emph{The farmers tie the vines to the elm trees.}
\end{quote}

Second declension but feminine: \emph{ulmus, ulmī f.} Italian farmers trained vines up elms. English \emph{elm} is a Germanic cousin.

\begin{cousinweb}
\emph{ulmus} is inherited from PIE \emph{h\textsubscript{1}élem} 'mountain elm' (EtymDB); English \emph{elm} is a Germanic cousin from the same root.
\end{cousinweb}

\begin{grammarlens}[title={What you've built}]
One more word for telling, explaining and describing at B1.
\end{grammarlens}

\subsection*{Your turn}
\begin{itemize}
\raggedright
  \item \emph{Say it:} \emph{ulmus, ulmī}
  \item \emph{Say it:} \emph{ulmus, ulmī}, once more
  \item \emph{Say it:} a sentence of your own with \emph{ulmus, ulmī}
  \item \emph{Recall:} say the Latin for a beech tree, then the Latin for a fir tree, then say \emph{ulmus, ulmī} again
\end{itemize}

\begin{tcolorbox}[breakable,colback=teal!4,colframe=teal!35!black,title={Before you move on}]
What does ulmus, ulmī mean? (\textquotedblleft{}An elm tree\textquotedblright{}.) Say it once more.
\end{tcolorbox}
\section[fraxinus, fraxinī]{fraxinus, fraxinī — an ash tree; an ash-wood spear}

\label{lesson:LA-C484-fraxinus}



\begin{quote}
Before the new one: say the Latin for a fir tree, then the Latin for an elm tree.
\end{quote}

\subsection*{You'll want to know: fraxinus, fraxinī}
\textbf{fraxinus, fraxinī} — \textquotedblleft{}an ash tree; an ash-wood spear\textquotedblright{}.

\begin{quote}
Ex fraxinō hastae fīunt.
\end{quote}

\begin{quote}
\emph{Spears are made from ash wood.}
\end{quote}

Second declension but feminine: \emph{fraxinus, fraxinī f.} Its hard, springy wood made spear shafts.

\begin{cousinweb}
\emph{fraxinus} is inherited from Proto-Italic \emph{fraksinos}, from PIE \emph{b\textsuperscript{h}rHǵ-s-inos} 'ash tree' (EtymDB).
\end{cousinweb}

\begin{grammarlens}[title={What you've built}]
One more word for telling, explaining and describing at B1.
\end{grammarlens}

\subsection*{Your turn}
\begin{itemize}
\raggedright
  \item \emph{Say it:} \emph{fraxinus, fraxinī}
  \item \emph{Say it:} \emph{fraxinus, fraxinī}, once more
  \item \emph{Say it:} a sentence of your own with \emph{fraxinus, fraxinī}
  \item \emph{Recall:} say the Latin for a fir tree, then the Latin for an elm tree, then say \emph{fraxinus, fraxinī} again
\end{itemize}

\begin{tcolorbox}[breakable,colback=teal!4,colframe=teal!35!black,title={Before you move on}]
What does fraxinus, fraxinī mean? (\textquotedblleft{}An ash tree; an ash-wood spear\textquotedblright{}.) Say it once more.
\end{tcolorbox}
\section[opācus, opāca, opācum]{opācus, opāca, opācum — shady, shaded; dark}

\label{lesson:LA-C484-opacus}



\begin{quote}
Before the new one: say the Latin for an elm tree, then the Latin for an ash tree.
\end{quote}

\subsection*{You'll want to know: opācus, opāca, opācum}
\textbf{opācus, opāca, opācum} — \textquotedblleft{}shady, shaded; dark\textquotedblright{}.

\begin{quote}
Aestāte viātōrēs locum opācum quaerunt.
\end{quote}

\begin{quote}
\emph{In summer travellers look for a shady place.}
\end{quote}

An adjective of three endings. English \emph{opaque} comes from it.

\begin{cousinweb}
The origin of \emph{opācus} is not settled: Lewis \& Short give no etymology and EtymDB gives no parent. English \emph{opaque} comes from it through French.
\end{cousinweb}

\begin{grammarlens}[title={What you've built}]
One more word for telling, explaining and describing at B1.
\end{grammarlens}

\subsection*{Your turn}
\begin{itemize}
\raggedright
  \item \emph{Say it:} \emph{opācus, opāca, opācum}
  \item \emph{Say it:} \emph{opācus, opāca, opācum}, once more
  \item \emph{Say it:} a sentence of your own with \emph{opācus, opāca, opācum}
  \item \emph{Recall:} say the Latin for an elm tree, then the Latin for an ash tree, then say \emph{opācus, opāca, opācum} again
\end{itemize}

\begin{tcolorbox}[breakable,colback=teal!4,colframe=teal!35!black,title={Before you move on}]
What does opācus, opāca, opācum mean? (\textquotedblleft{}Shady, shaded; dark\textquotedblright{}.) Say it once more.
\end{tcolorbox}
